\section{Design Exploration and Abandoned Routes}

This appendix is the honest record of routes tried and dropped. Keeping the dead ends visible is part of the integrity of the program. What follows collects the explorations that did not enter the committed model, the routes that violated a base commitment, and the design-option maps where nothing was ever committed. None of this is the theory. It is the search around the theory, kept so that the path is legible and the final conclusion is earned rather than asserted.

\subsection{The question behind all of it}

Every route in this appendix asks one question in a different dress. Can a confined body radiate.

The committed law, the knit, is collide then stream. Its collision phase reflects, and that reflection is what seals a body in place. A reflecting collision pins identity. The same reflection, though, blocks the only channel by which a disturbance could leave. So a body holds together but cannot shed what hits it. A self, the persistent self-maintaining knot the theory cares about, needs to shed disturbances. Shedding is its agency. A body that cannot radiate cannot settle, and a thing that cannot settle cannot keep an identity across many beats.

\begin{principle}[The radiation question]
The committed knit pins a body by reflecting its collisions. The same reflection seals the radiation channel. A self must shed disturbances to hold its identity. The explorations below all ask the same thing. What is the minimal addition, kept fully discrete, that lets a pinned body emit a disturbance and settle back to its ground.
\end{principle}

Each route is one answer to that question. Each was either ruled out for breaking a commitment or left uncommitted as a recorded option. The throughline at the end shows that all three corrected to the same single picture.

\subsection{Route one, the second-field gauge, abandoned}

The most fully worked route added a \textbf{second ternary field}, a discrete photon living on the links, alongside the matter tone carried by the sites of each dock. This is the discrete analogue of the continuous gauge field of textbook electromagnetism.

\subsubsection{The encoding}

The encoding is clean, and that cleanness is what made the route tempting. The ternary tone $\{-1, 0, +1\}$ is exactly the group of integers modulo three, with $-1$ standing for two. Addition is modulo three. Table~\ref{tab:z3-encoding} gives the dictionary.

\begin{table}[ht]
\centering
\begin{tabular}{@{}lccc@{}}
\toprule
as tone & $-1$ & $0$ & $+1$ \\
\midrule
as the modular group & $2$ & $0$ & $1$ \\
\bottomrule
\end{tabular}
\caption{The ternary tone read as the group of integers modulo three. The tone alphabet is itself the gauge group.}
\label{tab:z3-encoding}
\end{table}

So the discrete gauge group is the integers modulo three, the finite analogue of the continuous circle phase of ordinary electromagnetism. The continuous phase $e^{i\theta}$ becomes a single modular element. No real number appears anywhere. This is the whole appeal of a discrete gauge field on this substrate. The tone alphabet \emph{is} the gauge group, so the photon needs no new alphabet, only a new place to live.

\subsubsection{The two fields}

Under this route there are two fields per dock, both ternary, both discrete, as set out in Table~\ref{tab:two-fields}.

\begin{table}[ht]
\centering
\begin{tabular}{@{}llll@{}}
\toprule
field & lives on & values & meaning \\
\midrule
matter & each site of a dock & $\{-1,0,+1\}$ & the charged matter tone, the body \\
photon & each directed link & $\{-1,0,+1\}$ & the neutral quantum on that link \\
\bottomrule
\end{tabular}
\caption{The two fields of the abandoned gauge route. The matter charge of a dock is the signed sum of its site tones. The photon carries no matter charge, that neutrality is its whole point.}
\label{tab:two-fields}
\end{table}

A directed link and its reverse are the same undirected link, so the reverse link carries the negated value modulo three. There is one independent value per undirected link. The photon streams along its link to the next collinear link, one step per beat, exactly as matter streams but in the second field. This gives a massless propagating photon that reaches the wake, the open frontier, and is absorbed there.

\subsubsection{The coupling}

The coupling is one reversible move at each dock. An excited matter charge sitting in a fast site relaxes to a slow site and deposits a photon quantum on a link. The inverse of that move is absorption. The move acts on a triple, a fast site $L$, a slow site $S$, and a photon link $P$, chosen so that momentum is conserved, the displacement of $L$ equalling the displacement of $S$ plus the displacement of $P$. Table~\ref{tab:emission} is the exact permutation, a set of two-cycles, so it is its own inverse, an involution.

\begin{table}[ht]
\centering
\begin{tabular}{@{}llll@{}}
\toprule
before $(m_L, m_S, a_P)$ & after & move \\
\midrule
$(+1, 0, 0)$ & $(0, +1, +1)$ & emit \\
$(0, +1, +1)$ & $(+1, 0, 0)$ & absorb \\
$(-1, 0, 0)$ & $(0, -1, -1)$ & emit \\
$(0, -1, -1)$ & $(-1, 0, 0)$ & absorb \\
the other $23$ states & unchanged & inert \\
\bottomrule
\end{tabular}
\caption{The emission coupling as an exact involution on the local triple. An excited matter charge in the fast site drops to the slow site and deposits a photon on the link, and the reverse move re-excites it.}
\label{tab:emission}
\end{table}

On paper the tables are exactly conservative and reversible. Matter charge is conserved, because the photon is neutral and carries none of it. Momentum is conserved by the chosen triple. The move is a set of two-cycles plus fixed points, so applying it twice is the identity. The full beat becomes a fixed sequence of reversible sub-steps, matter collide, matter stream, the emission coupling, photon stream, with the only irreversibility living at the open boundary where the wake absorbs what falls behind it.

\begin{lemma}[The gauge route is conservative on paper]
The emission coupling conserves matter charge, because the photon is neutral. It conserves momentum, because the triple is chosen so the fast displacement equals the slow displacement plus the photon displacement. It is reversible, because it is an involution. Charge conservation together with the coupling yields a discrete Gauss law, the divergence of the gauge field around a dock equalling the matter charge change there.
\end{lemma}

\subsubsection{Why it was abandoned}

The route works on paper and still it was dropped, for one reason. It violates the one-tone commitment.

\begin{result}[One tone, no second field]
The gauge route adds a \textbf{second ternary field}, a whole new state per dock and per link, where the model holds that there is only one tone. The result that seemed to force the second field, that no fixed subset of directions can be a clean photon sector, rules out a photon \emph{sector}, not a photon \emph{wave}. The correct picture keeps one tone, with the photon as a wave, a phonon, of that single field. No second field is needed.
\end{result}

The gauge route is kept only as the richer-than-necessary alternative. It is a full discrete gauge theory, with a Gauss law and Wilson loops available if ever wanted later, including the textbook covariant-hop form where matter picks up a link's value as a phase when it hops. That covariant form needs the matter to carry an internal phase index beyond its charge, which is yet more structure. It is not the minimal design, and it is not committed.

\subsection{Route two, the discrete-coupling map, nothing committed}

A separate exploration asked what \emph{minimal coupling} could even mean when the base has no real numbers.

In textbook physics a continuous coupling constant sets how strongly matter and the photon interact. The base here is discrete. The tone is ternary, and the knit is a finite permutation. So a real coupling constant is not a legitimate base ingredient. It can only ever be emergent, never put in at the base.

\begin{principle}[A coupling is a move, not a number]
With no real numbers at the base, the coupling becomes one more entry in the collision table, a reversible charge-conserving swap between a matter site and a photon site. Its strength is binary, the move is in the table or it is not. This is exactly how the rest of the knit already works. The create move, the hop, the inert state are all discrete table entries, not numbers.
\end{principle}

If a tunable interaction is ever wanted, it stays discrete. One varies the \emph{set} of allowed emission triples, or the size of the mesh, never a real constant. This matches the methodology of the whole program, where robustness comes from size, never from a real parameter.

Concretely, the emission move is an excited fast matter charge dropping to a slower state and emitting a photon, with the photon carrying the momentum difference so the vector sum is exact. This was worked out as a clean reversible family of table entries over the allowed triples. Nothing here was committed. It is the concrete shape of a discrete coupling, recorded for discussion.

\subsection{Route three, the discrete-self map, nothing committed}

The broadest exploration laid out every option for making a confined body radiate, across four coupled problems, and ranked the options by how far each departs from the committed base. Nothing was committed. The value of the map is that it pins down exactly which options are alive and which are dead.

\subsubsection{A correction made along the way}

An early claim said emission needs multi-speed sites, on the grounds that one direction cannot split into two of the same length. That is wrong.

Momentum is a vector. Single-speed sites can conserve it by \textbf{recoil}, the matter changing direction rather than speed and emitting a photon to balance the books. Such triples genuinely exist among the 24 sites. So multi-speed sites are not strictly required. The one thing single-speed recoil does not conserve is a quadratic energy-like quantity. But that is not a base conservation law. Charge and momentum are the base conserved quantities, and recoil respects both.

\subsubsection{The four problems}

\begin{table}[ht]
\centering
\begin{tabular}{@{}p{0.24\textwidth}p{0.30\textwidth}p{0.36\textwidth}@{}}
\toprule
problem & what it asks & where it landed \\
\midrule
emission kinematics & do momentum-conserving emission triples exist & yes, abundant, not the obstacle. 192 triples on the single-speed sites, 336 on a two-speed set \\
sector assignment & which sites are matter, which are the neutral photon & the decisive gate, see below \\
reversibility and conservation & does the emission family form a clean reversible permutation & a finite check, several clean structures available \\
does it settle & does the emitted photon reach the wake and the body return to ground & open, empirical, the make-or-break, not yet confirmed \\
\bottomrule
\end{tabular}
\caption{The four coupled problems of the discrete-self map. Three are settled on paper. The fourth, whether a perturbed body actually radiates and settles, is the decisive unverified experiment.}
\label{tab:four-problems}
\end{table}

\subsubsection{The decisive result, and the wrong turn it caused}

The emission triples are abundant, so kinematics is not the obstacle. The real gate is that the photon must be \textbf{charge-neutral}, or emission would not conserve matter charge. A single ternary tone in a site carries charge, so the photon cannot be just any site. The question became whether there is a partition of the 24 sites into a neutral set and a charged set such that every emission triple has its emitted photon in the neutral set and its matter in the charged set.

\begin{proposition}[No clean photon sector among the sites]
Enumerated over the 24 sites, the coloring does not exist. Every one of the 24 sites appears as an emitted photon in some triple and as matter in another. No site can be consistently colored. So there is no clean photon \emph{sector} inside the directions.
\end{proposition}

This drove the wrong turn. The conclusion drawn at the time was that radiation needs a separate neutral field, and that inference led straight to the gauge route of the first subsection. The inference was mistaken.

\begin{result}[Sector ruled out, wave survives]
The statement no fixed direction-sector is the photon does not imply a separate field is the photon. It only rules out the photon being a static subset of sites. The photon can instead be a wave of the one tone field, which is not a sector and so is not touched by the no-coloring result. Matter is the tone bound, the photon is the tone waving, a neutral phonon. That is exactly mass-energy equivalence, and the single-field picture survives.
\end{result}

\subsubsection{The ranking, recorded but superseded}

The map ranked candidate site-sets by their distance from the committed base. The order ran from the single-speed 24 sites with recoil emission as the smallest change, then the 24 sites with one added still direction for a true zero-speed ground state, then a two-speed set where emission also conserves an energy-like quantity, then the full symmetry set whose half-integer directions break the integer lattice, then a separate tone field for the photon as the largest change. Three options sit out of bounds entirely, each breaking a commitment, as Table~\ref{tab:out-of-bounds} records.

\begin{table}[ht]
\centering
\begin{tabular}{@{}ll@{}}
\toprule
out-of-bounds option & commitment it breaks \\
\midrule
a real coupling constant & discreteness \\
a dissipative base knit & reversibility \\
non-local links & locality \\
\bottomrule
\end{tabular}
\caption{The three options ruled out at the base for breaking a standing commitment.}
\label{tab:out-of-bounds}
\end{table}

This whole ranking was framed around adding a second field or multi-speed sites, and the single-field phonon picture supersedes it. It is kept here as the map of the territory, not as a recommendation. A still home direction, a true rest slot, was among the options and is recorded as optional. It was later confirmed not to be required for binding, so a dock holds exactly 24 sites and no 25th.

\subsection{The throughline}

Every route in this appendix was an attempt to give a pinned body the agency to radiate without breaking a base commitment.

The gauge route works on paper but adds a forbidden second tone. The discrete-coupling map showed that a real coupling constant is not a base ingredient and that a discrete move is the legitimate form. The discrete-self map proved that emission kinematics is easy and that the photon cannot be a fixed sector, which was then misread as needing a second field.

\begin{result}[The earned conclusion]
The honest correction in all three routes is the same. Keep one tone. The photon is the tone waving, not a new field. These routes are the dead ends that made that clear, and they are kept so the conclusion is earned, not asserted.
\end{result}

The one thing still to measure, across all three routes, is the same open question. Run forward on the open mesh, does a perturbed bound charge genuinely radiate its disturbance to the wake and settle back to ground. That is the decisive experiment, and it remains the make-or-break for matter that holds an identity.
